\chapter{Eigenvalue Problems and Spectral Transforms}
\label{ch:eigenvalues}

\chapterorient{The eigenmode analysis asks a question none of the linear
solvers of this part can answer: not ``what field does this drive produce?''
but ``what fields can ring with no drive at all?'' Mathematically it asks for a
few eigenpairs of the generalized eigenproblem $\Kop\vx=\lambda\Mop\vx$---and,
maddeningly, it wants the \emph{smallest} few, the lowest resonances buried at
the bottom of an enormous spectrum. Plain iteration finds the \emph{largest}
eigenvalues, the exact opposite of what we need. This chapter develops the one
idea that turns the hard interior problem into an easy dominant one: the
\emph{spectral transform}. We watch power iteration climb toward the dominant
eigenvector, see why it climbs the wrong direction, and then bend the spectrum
with shift-and-invert so the eigenvalues we want become the ones iteration
naturally finds---paying for it with an inner linear solve. By the end you will
know why the eigenmode driver assembles its operators, hands them to a library,
and reads back converged modes, with no Palace-authored loop in sight.}

\section{The problem: a few low resonances of a huge pencil}
\label{sec:evp-problem}

The eigenmode analysis of \cref{ch:targets} seals a region of space with
conducting walls and asks what fields it can sustain on its own. We derived its
governing equation in \cref{ch:reductions-pde}: the source-free time-harmonic
curl--curl equation, discretized, is the \emph{generalized eigenvalue problem}
\begin{equation}
  \boxed{\;\Kop\,\vx \;=\; \lambda\,\Mop\,\vx,\qquad \lambda=\omega^2\;}
  \label{eq:evp-gevp}
\end{equation}
where $\Kop$ is the $\nu$-weighted curl--curl \emph{stiffness} and $\Mop$ the
$\varepsilon$-weighted \emph{mass} (\cref{ch:reductions-pde}). The pair
$(\Kop,\Mop)$ is called a \emph{matrix pencil}: not one matrix but two, and the
eigenproblem lives in the family $\Kop-\lambda\Mop$ of their combinations. Each
eigenvalue $\lambda$ gives a resonant frequency $\omega=\sqrt\lambda$; each
eigenvector $\vx$ gives the mode field. The word ``generalized'' refers to the
$\Mop$ on the right: an \emph{ordinary} eigenproblem reads $\mat{A}\vx=\lambda\vx$
with the identity in $\Mop$'s place, but here the $\varepsilon$-weighting puts a
genuine mass matrix there.

\begin{notationbox}
A \emph{pencil} is an ordered pair of matrices $(\Kop,\Mop)$ of the same size; we
write the generalized eigenproblem as $\Kop\vx=\lambda\Mop\vx$ or, equivalently,
$(\Kop-\lambda\Mop)\vx=\mathbf 0$. A scalar $\lambda$ for which $\Kop-\lambda\Mop$
is singular is an \emph{eigenvalue}; a nonzero $\vx$ in its null space is the
corresponding \emph{eigenvector}; together they form an \emph{eigenpair}
$(\lambda,\vx)$. When $\Kop$ and $\Mop$ are symmetric and $\Mop$ is positive
definite---as they are for a lossless cavity---all eigenvalues are real and the
eigenvectors can be chosen $\Mop$-orthogonal, $\inner{\Mop\vx_i}{\vx_j}=0$ for
$i\neq j$.
\end{notationbox}

So far this is bookkeeping. The difficulty---the one that shapes the entire
chapter---is \emph{which} eigenpairs we want, and \emph{how many}.

\paragraph{We want a few, not all.} A finely meshed three-dimensional cavity has
$N$ degrees of freedom running into the millions (\cref{ch:krylov}), so the
pencil has millions of eigenvalues. The engineer does not want them all. A
filter designer wants the lowest handful of resonances---the fundamental mode and
a few overtones, perhaps ten eigenpairs out of ten million. Computing the entire
spectrum would be absurd: it would cost more than a million linear solves and
produce a flood of numbers no one will read. The eigenmode analysis is, by
construction, a request for a \emph{small slice} of an enormous spectrum.

\paragraph{We want the smallest, which is the hard end.} Worse, the slice we want
sits at the \emph{bottom}. The lowest resonant frequencies are the smallest
$\lambda$, and the smallest nonzero eigenvalues of a matrix are, as we are about
to see, exactly the ones simple iteration is worst at finding. Iteration
naturally amplifies the \emph{largest} eigenvalues; the smallest are buried,
their signal swamped. Finding the interior or lower end of a spectrum is a
genuinely harder problem than finding the dominant end, and the rest of this
chapter is the story of how we cheat our way around that difficulty.

\Cref{fig:evp-setup} draws the situation: the pencil on the left, the vast
spectrum on the right, and the small wanted band marked at its low end.

\begin{figure}[t]
\centering
\begin{tikzpicture}[node distance=10mm]
  % the pencil, left
  \node[stage, minimum width=20mm, minimum height=10mm] (K) at (0,0.9) {$\Kop$};
  \node[stage, minimum width=20mm, minimum height=10mm] (M) at (0,-0.5) {$\Mop$};
  \node[font=\scriptsize\itshape, simGray, above=1mm of K] {stiffness};
  \node[font=\scriptsize\itshape, simGray, below=1mm of M] {mass};
  \node[draw=simGray, densely dashed, rounded corners=3pt, fit=(K)(M),
        inner sep=4pt, label={[font=\scriptsize\itshape,simGray]left:pencil}] (pen) {};
  % the eigenproblem arrow
  \node[font=\small] (eq) at (3.0,0.2) {$\Kop\vx=\lambda\Mop\vx$};
  \draw[flow] (pen.east) -- (eq.west);
  % the spectrum axis, right
  \draw[flow, simGray] (4.7,-1.2) -- (12.6,-1.2)
    node[right, font=\scriptsize\itshape, text=simGray] {$\lambda$};
  % many eigenvalue ticks, dense toward the right (large)
  \foreach \x in {5.1,5.45,5.95,6.6,7.4,8.3,9.2,10.0,10.6,11.0,11.3,11.55,11.75,11.9,12.05,12.18,12.3}{
    \draw[simBlue, thick] (\x,-1.32) -- (\x,-1.08);}
  % the wanted low band
  \fill[simRust!18] (4.95,-1.55) rectangle (6.05,-0.85);
  \draw[simRust, thick] (4.95,-1.55) rectangle (6.05,-0.85);
  \node[font=\scriptsize\bfseries, simRust, align=center] at (5.5,-0.55)
    {wanted:\\lowest few};
  % the "dominant" end iteration finds
  \node[font=\scriptsize\itshape, simGray, align=center] at (12.0,-0.55)
    {what iteration\\finds (dominant)};
  \draw[refedge, simGray] (12.0,-0.85) -- (12.2,-1.06);
\end{tikzpicture}
\caption[The generalized eigenproblem and its wanted low end]{The eigenmode
problem in one picture. The assembled pencil $(\Kop,\Mop)$ (left) defines the
generalized eigenproblem $\Kop\vx=\lambda\Mop\vx$, whose spectrum (right) holds
millions of eigenvalues. The engineer wants only the \emph{lowest few} (rust
band)---the fundamental resonance and a small number of overtones. But simple
iteration naturally surfaces the \emph{dominant} (largest) eigenvalues at the far
right, the exact opposite end. Bridging that mismatch is the business of the
chapter.}
\label{fig:evp-setup}
\end{figure}

\section{Power iteration: climbing toward the dominant direction}
\label{sec:power}

Before we can bend the spectrum we must understand the one iteration that all of
eigenvalue computation is built on. It is almost embarrassingly simple, and
understanding exactly what it does---and what it cannot do---is the whole key.

Take an ordinary eigenproblem $\mat{A}\vx=\mu\vx$ for the moment (we restore the
pencil shortly). \emph{Power iteration} does just one thing: it multiplies a
vector by $\mat{A}$ over and over, renormalizing to keep it bounded.
\begin{equation}
  \vx_{k+1} \;=\; \frac{\mat{A}\,\vx_k}{\norm{\mat{A}\,\vx_k}}.
  \label{eq:power-step}
\end{equation}
Why this converges to anything is the eigenvalue story in miniature. Expand the
starting vector in the eigenbasis $\{\vu_1,\dots,\vu_N\}$ of $\mat{A}$, ordered so
that $\abs{\mu_1}>\abs{\mu_2}\ge\dots\ge\abs{\mu_N}$:
\begin{equation}
  \vx_0 = c_1\vu_1 + c_2\vu_2 + \dots + c_N\vu_N .
\end{equation}
Applying $\mat{A}$ once multiplies each component $\vu_i$ by its eigenvalue
$\mu_i$; applying it $k$ times multiplies by $\mu_i^k$:
\begin{equation}
  \mat{A}^k\vx_0
    = c_1\mu_1^k\vu_1 + c_2\mu_2^k\vu_2 + \dots
    = \mu_1^k\Bigl(c_1\vu_1
        + c_2\Bigl(\tfrac{\mu_2}{\mu_1}\Bigr)^{\!k}\vu_2 + \dots\Bigr).
  \label{eq:power-expand}
\end{equation}
Because $\abs{\mu_2/\mu_1}<1$, every ratio $(\mu_i/\mu_1)^k$ decays to zero as
$k$ grows. The $\vu_1$ term, scaled by the largest eigenvalue, swamps all the
others. After enough steps the iterate points (up to sign and scale) in the
direction of $\vu_1$, the \emph{dominant eigenvector}---the eigenvector belonging
to the eigenvalue of largest magnitude. \Cref{fig:power-rotate} shows the
geometry: a generic starting vector swings, step by step, toward the dominant
axis as the subdominant component shrinks.

\begin{figure}[t]
\centering
\begin{tikzpicture}[scale=1.0]
  % axes: dominant eigenvector u1 (horizontal-ish), subdominant u2 (vertical-ish)
  \draw[->, simGray] (0,0) -- (5.0,0) node[right,font=\scriptsize]{$\vu_1$ (dominant)};
  \draw[->, simGray] (0,0) -- (0,3.4) node[above,font=\scriptsize]{$\vu_2$};
  % a sequence of normalized vectors rotating toward u1
  \foreach \ang/\op in {65/0.30, 46/0.42, 30/0.54, 18/0.66, 10/0.80, 5/1.0}{
    \draw[-{Stealth[length=2.2mm]}, very thick, simBlue, opacity=\op]
      (0,0) -- (\ang:3.0);
  }
  % final converged vector emphasized
  \draw[-{Stealth[length=2.6mm]}, very thick, simRust] (0,0) -- (5:3.0);
  \node[font=\scriptsize, simRust] at (5:3.4) {$\vx_k$};
  \node[font=\scriptsize, simBlue] at (66:3.3) {$\vx_0$};
  % arc showing the swing
  \draw[->, simGray, densely dotted] (62:2.0) arc (62:9:2.0);
  \node[font=\scriptsize\itshape, simGray] at (35:2.35) {each step};
  \node[font=\scriptsize\itshape, simGray] at (33:1.85) {$\times\,\mu_2/\mu_1$};
\end{tikzpicture}
\caption[Power iteration rotates toward the dominant eigenvector]{Power
iteration as a rotation. The starting vector $\vx_0$ (faint blue) has a component
along every eigenvector; each multiplication by $\mat{A}$ multiplies the
$\vu_2$ component by $\mu_2/\mu_1<1$ relative to the $\vu_1$ component, so the
iterate swings, step by step, toward the dominant axis $\vu_1$. The converged
$\vx_k$ (rust) points along $\vu_1$. The speed of the swing is set by the ratio
$\abs{\mu_2/\mu_1}$: the closer to $1$, the slower the rotation.}
\label{fig:power-rotate}
\end{figure}

\subsection{The Rayleigh quotient: reading off the eigenvalue}

Power iteration delivers a \emph{direction}---the eigenvector. To get the
\emph{eigenvalue} we evaluate the \emph{Rayleigh quotient}, the single most
useful scalar in eigenvalue computation:
\begin{equation}
  \rho(\vx) \;=\; \frac{\inner{\mat{A}\vx}{\vx}}{\inner{\vx}{\vx}}.
  \label{eq:rayleigh}
\end{equation}
If $\vx$ is exactly an eigenvector, $\mat{A}\vx=\mu\vx$, then
$\rho(\vx)=\inner{\mu\vx}{\vx}/\inner{\vx}{\vx}=\mu$: the quotient returns the
eigenvalue exactly. If $\vx$ is only \emph{near} an eigenvector, $\rho(\vx)$ is a
remarkably good estimate---for a symmetric $\mat{A}$ it is accurate to
\emph{second} order in the eigenvector error, so a direction good to three digits
yields an eigenvalue good to six. As power iteration rotates the iterate toward
$\vu_1$, the Rayleigh quotient $\rho(\vx_k)$ closes in on $\mu_1$, and the two
together---direction from \cref{eq:power-step}, value from
\cref{eq:rayleigh}---form a complete eigenpair estimate.

\begin{figure}[t]
\centering
\begin{tikzpicture}[node distance=8mm]
  \node[opbox, minimum width=20mm] (x) {$\vx$};
  \node[opbox, minimum width=22mm, right=12mm of x] (Ax) {$\mat{A}\vx$};
  \node[opbox, fill=simBgGold, draw=simGold, minimum width=30mm, right=12mm of Ax] (num)
    {$\inner{\mat{A}\vx}{\vx}$};
  \node[opbox, fill=simBgGold, draw=simGold, minimum width=26mm, below=6mm of num] (den)
    {$\inner{\vx}{\vx}$};
  \node[product, minimum width=20mm, right=14mm of num, yshift=-7mm] (rho)
    {$\rho(\vx)$};
  \draw[flow] (x) -- node[above,font=\scriptsize]{apply} (Ax);
  \draw[flow] (Ax) -- node[above,font=\scriptsize]{dot} (num);
  \draw[flow] (x.south) to[out=-90,in=180] (den.west);
  \draw[flow] (num.east) -- node[above,font=\scriptsize]{$\div$} (rho.north west);
  \draw[flow] (den.east) -- (rho.south west);
  \node[font=\scriptsize\itshape, simGray, below=1mm of den]
    {if $\mat{A}\vx=\mu\vx$ then $\rho=\mu$ exactly};
\end{tikzpicture}
\caption[The Rayleigh quotient]{The Rayleigh quotient $\rho(\vx)=\inner{\mat{A}
\vx}{\vx}/\inner{\vx}{\vx}$ turns a vector into an eigenvalue estimate with one
operator apply and two inner products. On an exact eigenvector it returns the
eigenvalue exactly; near an eigenvector it is accurate to second order in the
direction error. Paired with power iteration's direction estimate, it completes
the eigenpair.}
\label{fig:rayleigh}
\end{figure}

\subsection{Several at once: subspace iteration}

We want a \emph{handful} of eigenpairs, not one. The natural generalization
iterates a whole \emph{block} of vectors rather than a single one.
\emph{Subspace} (or \emph{orthogonal}) iteration carries an $N\times p$ matrix
$\mat{X}_k$ of $p$ columns, applies $\mat{A}$ to all of them, and then
\emph{re-orthonormalizes} the columns:
\begin{equation}
  \mat{Y} = \mat{A}\,\mat{X}_k,\qquad
  \mat{X}_{k+1} = \operatorname{orth}(\mat{Y}),
\end{equation}
where $\operatorname{orth}$ is a QR-style orthonormalization
(\cref{ch:krylovschur}). Without the re-orthonormalization every column would
collapse onto the same dominant eigenvector; the orthogonalization forces the
columns apart, so they converge to the $p$ \emph{dominant} eigenvectors
$\vu_1,\dots,\vu_p$ simultaneously. The Rayleigh quotient generalizes too: the
small $p\times p$ matrix $\mat{X}_k^{\!\top}\mat{A}\,\mat{X}_k$ has eigenvalues
that approximate the $p$ dominant eigenvalues---the \emph{Rayleigh--Ritz}
projection that the Krylov methods of \cref{ch:krylovschur} sharpen into their
final form.

\begin{keyidea}
Power iteration, and its block form subspace iteration, converge to the
\emph{dominant} eigenpairs---those of \emph{largest} magnitude. The Rayleigh
quotient $\rho(\vx)=\inner{\mat{A}\vx}{\vx}/\inner{\vx}{\vx}$ reads the eigenvalue
off the converged direction. The convergence rate is governed by the
\emph{ratio} of consecutive eigenvalues $\abs{\mu_2/\mu_1}$: well-separated
eigenvalues converge fast, clustered ones agonizingly slowly. This is the engine
under every eigensolver---but on its own it finds exactly the wrong end of the
spectrum for the eigenmode problem.
\end{keyidea}

\subsection{Why this is not enough}

Two facts about power iteration doom it as a direct attack on the eigenmode
problem, and naming them sharply is what motivates the spectral transform.

\emph{First, it finds the largest, and we want the smallest.} The eigenmode
analysis wants the lowest resonances, the smallest nonzero $\lambda$. Power
iteration delivers the dominant eigenvalue, $\lambda_{\max}$---the
\emph{highest} resonance, of no interest, and in fact a meaningless mesh-resolution
artifact at the top of the discrete spectrum.

\emph{Second, it is slow exactly where we need it.} The convergence rate is
$\abs{\mu_2/\mu_1}$. The low resonances of a cavity are typically
\emph{clustered}---a fundamental and its near-degenerate overtones sit close
together---so consecutive eigenvalues have ratios near $1$, and power iteration
crawls. The very structure of the problem we care about is the structure power
iteration handles worst.

We could try to fix the first problem by iterating with $\mat{A}^{-1}$ instead of
$\mat{A}$: the eigenvalues of $\mat{A}^{-1}$ are $1/\mu_i$, so the smallest $\mu_i$
becomes the largest $1/\mu_i$, and power iteration on $\mat{A}^{-1}$ converges to
the \emph{smallest} eigenvector. This is \emph{inverse iteration}, and it is the
right instinct---but it points only at the very smallest eigenvalue, does nothing
for the clustering, and gives us no control over \emph{where} in the spectrum to
look. The spectral transform is inverse iteration grown up: it lets us aim.

\section{The spectral transform: shift-and-invert}
\label{sec:shift-invert}

Here is the central idea of the chapter. We cannot easily find interior or low
eigenvalues directly. So we \emph{transform} the operator into a new one whose
\emph{dominant} eigenvalues are precisely the original eigenvalues we want---and
then we run the power/subspace iteration we already understand on the transformed
operator, where it works beautifully.

\subsection{Deriving the eigenvalue map}

Start from the generalized eigenproblem \cref{eq:evp-gevp},
$\Kop\vx=\lambda\Mop\vx$. Choose a \emph{shift} $\sigma$---a number near the
eigenvalues we are hunting---and rearrange. Subtract $\sigma\Mop\vx$ from both
sides:
\begin{equation}
  \Kop\vx - \sigma\Mop\vx = \lambda\Mop\vx - \sigma\Mop\vx
  \quad\Longrightarrow\quad
  (\Kop-\sigma\Mop)\,\vx = (\lambda-\sigma)\,\Mop\,\vx .
\end{equation}
Now suppose $\lambda\neq\sigma$, so $\Kop-\sigma\Mop$ is invertible on $\vx$.
Multiply both sides by $(\Kop-\sigma\Mop)^{-1}$ and divide by the scalar
$(\lambda-\sigma)$:
\begin{equation}
  \frac{1}{\lambda-\sigma}\,\vx
    \;=\; (\Kop-\sigma\Mop)^{-1}\,\Mop\,\vx .
\end{equation}
Read this as an \emph{ordinary} eigenproblem for the transformed operator
\begin{equation}
  \boxed{\;\mat{T}_\sigma \;=\; (\Kop-\sigma\Mop)^{-1}\,\Mop,
    \qquad \mat{T}_\sigma\,\vx \;=\; \theta\,\vx,
    \quad \theta=\frac{1}{\lambda-\sigma}\;}
  \label{eq:shift-invert}
\end{equation}
The eigenvectors are \emph{unchanged}---$\vx$ is an eigenvector of $\mat{T}_\sigma$
exactly when it is an eigenvector of the pencil---but the eigenvalues are mapped
by
\begin{equation}
  \lambda \;\longmapsto\; \theta = \frac{1}{\lambda-\sigma}.
  \label{eq:eigmap}
\end{equation}

Now examine what this map does. The transformed eigenvalue $\theta=1/(\lambda
-\sigma)$ is \emph{large in magnitude exactly when $\lambda$ is close to
$\sigma$}, and small when $\lambda$ is far from $\sigma$. An eigenvalue sitting
right at the shift is sent to infinity; one far away is sent near zero. So the
eigenvalues of the original problem \emph{nearest the shift $\sigma$} become the
\emph{dominant} eigenvalues of $\mat{T}_\sigma$---and those are exactly the ones
power and subspace iteration find first and fastest.

\begin{keyidea}
\textbf{Shift-and-invert turns a hard interior eigenvalue problem into an easy
dominant one.} Iterating against $\mat{T}_\sigma=(\Kop-\sigma\Mop)^{-1}\Mop$
instead of the pencil maps each eigenvalue $\lambda\mapsto\theta=1/(\lambda
-\sigma)$, so the eigenvalues \emph{nearest the shift} $\sigma$ become the
\emph{largest} (dominant) eigenvalues of $\mat{T}_\sigma$. ``Find the eigenvalues
near $\sigma$''---hard, interior, what the eigenmode problem actually
needs---becomes ``find the dominant eigenvalues of $\mat{T}_\sigma$''---easy,
exactly what iteration does. The price is that every application of
$\mat{T}_\sigma$ requires an \emph{inner linear solve} of $(\Kop-\sigma\Mop)$,
turning each iteration step into a \code{ksp\_solve} (\cref{ch:krylov}).
\end{keyidea}

\subsection{The spectrum reordering, made vivid}

The map \cref{eq:eigmap} does more than enlarge the wanted eigenvalues: it
\emph{reorders} the entire spectrum. The next worked example makes the reordering
concrete with a four-eigenvalue spectrum, and \cref{fig:spectrum-reorder}---the
hero figure of the chapter---draws the before and after.

\begin{workedexample}{Shift-and-invert reorders a small spectrum}{reorder}
A pencil $(\Kop,\Mop)$ has four eigenvalues
\[
  \lambda_1 = 1,\quad \lambda_2 = 4,\quad \lambda_3 = 9,\quad \lambda_4 = 16,
\]
and we want the two nearest $\lambda=4$. Plain subspace iteration on the pencil
would converge to $\lambda_4=16$ and $\lambda_3=9$ first---the dominant pair, the
\emph{wrong} ones. Choose the shift $\sigma=4.5$, just above the target, and form
the transformed eigenvalues $\theta_i=1/(\lambda_i-\sigma)$:
\[
  \begin{aligned}
  \theta_1 = \frac{1}{1-4.5} = -0.286,\quad
  &\theta_2 = \frac{1}{4-4.5} = -2.0,\\
  \theta_3 = \frac{1}{9-4.5} = 0.222,\quad
  &\theta_4 = \frac{1}{16-4.5} = 0.087.
  \end{aligned}
\]
Rank them by magnitude:
\[
  \abs{\theta_2}=2.0 \;>\; \abs{\theta_1}=0.286 \;>\; \abs{\theta_3}=0.222
  \;>\; \abs{\theta_4}=0.087.
\]
The order has \emph{inverted}. In the original spectrum $\lambda_4=16$ was
dominant and $\lambda_2=4$ was a small interior value; after the transform
$\theta_2$ (from $\lambda_2=4$, nearest the shift) is by far the largest, and
$\theta_1$ (from $\lambda_1=1$, the next nearest) is second. Subspace iteration on
$\mat{T}_\sigma$ now converges \emph{first} to exactly the two eigenpairs we
wanted---those near $\sigma=4.5$---and the convergence is fast, because
$\abs{\theta_1/\theta_2}=0.286/2.0\approx 0.14$ is a small ratio, far from the
$\abs{\lambda_3/\lambda_4}=9/16\approx 0.56$ ratio that throttled the untransformed
iteration. We get the right eigenpairs \emph{and} faster convergence, both from
the single act of bending the spectrum.

\medskip
\emph{Reading the answer back.} The iteration returns $\theta_2=-2.0$. We recover
the physical eigenvalue by inverting \cref{eq:eigmap}:
$\lambda = \sigma + 1/\theta = 4.5 + 1/(-2.0) = 4.5 - 0.5 = 4$, exactly
$\lambda_2$. The eigenvector needs no untransforming---it is the \emph{same}
vector $\vx$ in both problems.
\end{workedexample}

\begin{figure}[t]
\centering
\begin{tikzpicture}
\begin{axis}[
    name=before,
    width=0.82\linewidth, height=34mm,
    xlabel={original eigenvalue $\lambda$},
    xmin=-1, xmax=18, ymin=0, ymax=1.4,
    axis lines=left, ytick=\empty, xtick={1,4,9,16},
    xticklabel style={font=\scriptsize}, xlabel style={font=\scriptsize},
    clip=false,
  ]
  % the four eigenvalues as stems
  \addplot[ycomb, very thick, simBlue, mark=*, mark size=2pt] coordinates {
    (1,1)(4,1)(9,1)(16,1)};
  % shift marker
  \draw[simRust, thick, densely dashed] (axis cs:4.5,0) -- (axis cs:4.5,1.3);
  \node[font=\scriptsize, simRust] at (axis cs:4.5,1.42) {shift $\sigma=4.5$};
  % wanted band
  \node[font=\scriptsize\itshape, simGray] at (axis cs:13,1.15) {dominant end};
  \node[font=\scriptsize, simBlue, anchor=south] at (axis cs:1,1) {$\lambda_1$};
  \node[font=\scriptsize, simBlue, anchor=south] at (axis cs:4,1) {$\lambda_2$};
  \node[font=\scriptsize, simBlue, anchor=south] at (axis cs:9,1) {$\lambda_3$};
  \node[font=\scriptsize, simBlue, anchor=south] at (axis cs:16,1) {$\lambda_4$};
\end{axis}
\begin{axis}[
    name=after, at={(before.south west)}, anchor=north west, yshift=-9mm,
    width=0.82\linewidth, height=40mm,
    xlabel={transformed eigenvalue $\theta=1/(\lambda-\sigma)$},
    ylabel={}, xmin=-2.6, xmax=2.6, ymin=0, ymax=2.4,
    axis lines=left, ytick=\empty,
    xtick={-2.0,-0.286,0.087,0.222}, xticklabels={$-2.0$,,,},
    xticklabel style={font=\scriptsize}, xlabel style={font=\scriptsize},
    clip=false,
  ]
  % transformed eigenvalues, height = magnitude (dominance)
  \addplot[ycomb, very thick, simRust, mark=*, mark size=2pt] coordinates {
    (-2.0,2.0)(-0.286,0.286)(0.222,0.222)(0.087,0.087)};
  \node[font=\scriptsize, simRust, anchor=south] at (axis cs:-2.0,2.0) {$\theta_2$ (was $\lambda_2$)};
  \node[font=\scriptsize, simGray, anchor=south] at (axis cs:-0.286,0.286) {$\theta_1$};
  \node[font=\scriptsize, simGray, anchor=south west] at (axis cs:0.222,0.222) {$\theta_3$};
  \node[font=\scriptsize, simGray, anchor=south west] at (axis cs:0.087,0.087) {$\theta_4$};
  \node[font=\scriptsize\itshape, simRust] at (axis cs:-1.4,1.7) {now dominant};
\end{axis}
\end{tikzpicture}
\caption[Shift-and-invert reorders the spectrum]{The spectral transform of
\cref{wex:reorder}, drawn. \emph{Top:} the original spectrum, four eigenvalues
$\{1,4,9,16\}$; the wanted pair sits near the shift $\sigma=4.5$ (rust dashed),
in the interior, while iteration's natural target is the dominant end at the
right. \emph{Bottom:} after the map $\theta=1/(\lambda-\sigma)$, the stem heights
show transformed magnitude (dominance). $\theta_2$ (from $\lambda_2=4$, nearest
the shift) towers over the rest---it has become the dominant eigenvalue---and
$\theta_1$ is second. Iteration on the transformed operator now finds exactly the
near-$\sigma$ pair first. Bending the spectrum turns the hard interior request
into the easy dominant one.}
\label{fig:spectrum-reorder}
\end{figure}

\subsection{The cost: an inner linear solve per apply}

The transform is not free, and the price is precisely located. To run iteration
on $\mat{T}_\sigma=(\Kop-\sigma\Mop)^{-1}\Mop$ we must, at every step, compute
$\mat{T}_\sigma\vv$ for the current iterate $\vv$. Written out, that is three
moves:
\begin{equation}
  \vv
  \;\xrightarrow{\ \text{apply }\Mop\ }\;
  \vw = \Mop\vv
  \;\xrightarrow{\ \text{inner solve}\ }\;
  \vy = (\Kop-\sigma\Mop)^{-1}\vw
  \;\xrightarrow{\ \text{untransform}\ }\;
  \vy .
  \label{eq:siapply}
\end{equation}
The first move is a cheap matrix--vector product against the mass matrix. The
\emph{second} move is the expensive one: applying $(\Kop-\sigma\Mop)^{-1}$ means
\emph{solving the linear system} $(\Kop-\sigma\Mop)\vy=\vw$. That is exactly the
sparse linear solve of \cref{ch:krylov}---a \code{ksp\_solve}---and because the
shifted matrix $\Kop-\sigma\Mop$ is large and sparse, it is itself solved
iteratively, almost always with a preconditioner (\cref{ch:multigrid}). The third
move is per-backend coordinate bookkeeping (an optional scaling un-twist) that
does not change the direction.

So the body of one shift-invert iteration step is the composition
\begin{equation}
  \text{apply }\Mop \;\rhd\; \text{inner }\code{ksp\_solve} \;\rhd\;
  \text{scale-untransform},
\end{equation}
and the \emph{outer} eigenvalue iteration wraps this body in its power/subspace
loop. \Cref{fig:si-apply} draws the body. This is the structural heart of the
shift-invert eigensolve: \emph{an outer eigenvalue iteration whose every step
contains an inner linear solve}. The solver of \cref{ch:krylov,ch:gmres} that we
built to solve $\mat{A}\vx=\vb$ now reappears as a \emph{subroutine called once
per eigen-iteration step}.

\begin{figure}[t]
\centering
\begin{tikzpicture}[node distance=9mm]
  \node[data, minimum width=14mm] (v) {$\vv$};
  \node[opbox, minimum width=22mm, right=10mm of v] (Mv) {apply $\Mop$};
  \node[extern, minimum width=34mm, right=10mm of Mv] (solve)
    {\textbf{inner }\code{ksp\_solve}\\[1pt]\footnotesize $(\Kop-\sigma\Mop)^{-1}$};
  \node[opbox, minimum width=26mm, right=10mm of solve] (un) {scale-untransform};
  \node[product, minimum width=20mm, right=10mm of un] (out) {$\mat{T}_\sigma\vv$};
  \draw[flow] (v) -- (Mv);
  \draw[flow] (Mv) -- node[above,font=\scriptsize]{$\vw$} (solve);
  \draw[flow] (solve) -- node[above,font=\scriptsize]{$\vy$} (un);
  \draw[flow] (un) -- (out);
  % annotations
  \node[font=\scriptsize\itshape, simGray, below=2mm of Mv] {cheap matvec};
  \node[font=\scriptsize\itshape, simRust, below=2mm of solve, align=center]
    {the cost:\\a full linear solve\\(\cref{ch:krylov,ch:multigrid})};
  \node[font=\scriptsize\itshape, simGray, below=2mm of un] {coordinate twist};
  % the outer loop wrapping
  \draw[backflow] (out.north) to[out=90,in=90]
    node[above,font=\scriptsize,simRust]{outer eigen-iteration repeats} (v.north);
\end{tikzpicture}
\caption[The shift-invert apply and its inner solve]{One application of the
transformed operator $\mat{T}_\sigma=(\Kop-\sigma\Mop)^{-1}\Mop$, the body of the
eigen-iteration. A matrix--vector product against the mass matrix $\Mop$ (cheap)
feeds an \emph{inner linear solve} of the shifted operator $\Kop-\sigma\Mop$
(expensive: a full \code{ksp\_solve}, \cref{ch:krylov}, usually preconditioned,
\cref{ch:multigrid}), followed by a per-backend scale-untransform. The outer
eigenvalue iteration (dashed rust) repeats this body each step. The linear solver
of the earlier chapters becomes a subroutine of the eigensolver.}
\label{fig:si-apply}
\end{figure}

\section{The eigensolver as a black box}
\label{sec:blackbox}

We now know, in principle, how to find the low resonances: transform the spectrum
with a shift, run subspace iteration on the transformed operator, recover the
eigenpairs. But Palace does \emph{not} write this loop. The eigenmode analysis is
the one analysis in the book whose ``solve'' stage is a single, opaque call to a
specialized library---and this section explains why, and what it means for the
shape of the code.

\subsection{The fourth solve corner}

\Cref{ch:situating} sorted every analysis's middle stage into one of \emph{four
solve corners}. Three of them are loops Palace writes: the fixed-operator family
solve (electrostatics), the operator-varying sweep (driven), the sequential fold
(transient). The fourth corner is different in kind: the \emph{opaque black-box
solve}, and the eigenmode analysis is its sole inhabitant. As \cref{ch:operators}
put it, a solver is an operator you \emph{apply}---and the eigensolver takes that
to its limit. The entire outer iteration lives \emph{inside a library}; Palace
hands over the operators and reads back the answer, authoring no loop at all.

The reason is practical and decisive. The outer eigenvalue iteration is not the
simple power iteration of \cref{sec:power}; production eigensolvers use
\emph{Krylov--Schur} or \emph{Arnoldi} methods (\cref{ch:krylovschur}), which
maintain and periodically ``restart'' a growing Krylov basis, run delicate
re-orthogonalization, apply locking and purging of converged pairs, and decide
convergence by subtle residual criteria. This machinery is intricate, easy to get
subtly wrong, and already implemented to a very high standard in specialized
libraries---SLEPc, ARPACK. Re-implementing it would be folly. So Palace composes
against it: it assembles $(\Kop,\Mop)$, configures the shift and the inner solver,
and calls the library's \code{Solve}.

\subsection{Reverse communication: the library calls back}

There is one wrinkle that makes the black box especially elegant. The library
drives the iteration, but it does not contain Palace's operators---it cannot,
because those operators are matrix-free and live in Palace's data structures
(\cref{ch:matrixfree}). So the library and Palace cooperate through
\emph{reverse communication}: the library runs its loop, and whenever it needs a
matrix--vector product, it \emph{pauses and calls back} into Palace, saying ``apply
the operator to this vector and hand me the result.'' Palace supplies the
\code{apply}---which, for shift-invert, is the whole inner-solve body of
\cref{fig:si-apply}---and the library resumes.

\Cref{fig:blackbox-rci} draws the control flow. The outer loop is the library's;
the apply callbacks are Palace's; and the line between them is the one interface
of \cref{ch:operators}: \code{apply(op, v)}. The eigensolver never sees Palace's
mesh, materials, or assembly; Palace never sees the library's Krylov basis,
restart logic, or convergence test. Each side holds what it owns and speaks to the
other only through the apply interface.

\begin{figure}[t]
\centering
\begin{tikzpicture}[node distance=10mm]
  % the library box (drives the loop)
  \node[extern, minimum width=40mm, minimum height=26mm] (lib) at (0,0) {};
  \node[font=\footnotesize\bfseries, simViolet] at (lib.north) [yshift=-3.5mm]
    {eigensolver library};
  \node[font=\scriptsize\itshape, simGray, align=center] at (lib.center) [yshift=2mm]
    {Krylov--Schur loop\\(\cref{ch:krylovschur}):\\restart, orthogonalize,\\
     lock, convergence test};
  \node[font=\scriptsize, simViolet] at (lib.south) [yshift=3.5mm] {SLEPc / ARPACK};
  % Palace apply box
  \node[stage, minimum width=34mm, minimum height=20mm, right=34mm of lib] (pal) {};
  \node[font=\footnotesize\bfseries, simBlue] at (pal.north) [yshift=-3.5mm]
    {Palace \code{apply}};
  \node[font=\scriptsize, simInk, align=center] at (pal.center) [yshift=-1mm]
    {$\vv\mapsto\mat{T}_\sigma\vv$:\\apply $\Mop$, inner solve,\\untransform
     (\cref{fig:si-apply})};
  % the reverse-communication callback arrows
  \draw[flow] (lib.east|-0,0.5) -- node[above,font=\scriptsize]{``apply to $\vv$''}
    (pal.west|-0,0.5);
  \draw[backflow] (pal.west|-0,-0.5) -- node[below,font=\scriptsize]{result $\mat{T}_\sigma\vv$}
    (lib.east|-0,-0.5);
  % setup in, eigenpairs out
  \node[data, align=center, text width=22mm, above=8mm of lib] (in)
    {assemble $(\Kop,\Mop)$,\\ shift $\sigma$, inner ksp};
  \draw[flow] (in) -- (lib);
  \node[product, align=center, text width=22mm, below=8mm of lib] (out)
    {converged\\eigenpairs $(\lambda_i,\vx_i)$};
  \draw[flow] (lib) -- (out);
  % labels for who owns what
  \node[font=\scriptsize\itshape, simViolet, above=0.5mm of in] {one opaque call};
  \node[font=\scriptsize\bfseries, simInk] at ($(lib.east)!0.5!(pal.west)+(0,1.6)$)
    {one interface: \code{apply}};
\end{tikzpicture}
\caption[The eigensolver as a black box with reverse-communication callbacks]{The
black-box eigensolve. Palace assembles the pencil $(\Kop,\Mop)$, picks a shift
$\sigma$ and an inner linear solver, and makes \emph{one opaque call} into a
specialized library (SLEPc, ARPACK). The library \emph{drives} the outer
iteration---the Krylov--Schur restart, orthogonalization, locking, and convergence
test of \cref{ch:krylovschur}---and whenever it needs a matrix--vector product it
\emph{calls back} into Palace through the one \code{apply} interface of
\cref{ch:operators} (\emph{reverse communication}). Palace's apply is the entire
shift-invert body of \cref{fig:si-apply}. Converged eigenpairs come back out. No
Palace-authored loop exists.}
\label{fig:blackbox-rci}
\end{figure}

\subsection{Contrast with the linear-solve corners}

The contrast with the other three solve corners is sharp and worth stating
plainly. In the electrostatic, driven, and transient analyses, Palace
\emph{writes the loop}: the per-terminal family loop, the frequency sweep, the
time march are all Palace source, and the linear solver they call is a
\code{ksp\_solve} \emph{whose iteration Palace also owns} (\cref{ch:krylov}).
There the solver is a constructed operator Palace builds and drives.

In the eigenmode analysis the ownership inverts. The \emph{outer} loop belongs to
the library; Palace owns only the \emph{apply} the library calls back for. The
eigensolve is therefore the deepest composition in the book: a library-owned outer
iteration, calling back into a Palace-owned shift-invert body, which itself
contains a Palace-owned \emph{inner} \code{ksp\_solve}, which itself calls a
preconditioner. Three nested layers of ``a solver is an operator''
(\cref{ch:operators}), and the outermost one is not even Palace's code.

\begin{pitfall}
Because the outer eigen-iteration is library-owned, the eigenmode analysis is the
one place in this book where the constructive ``iteration rotation'' of the later
chapters \emph{cannot reach the loop}: there is no Palace-authored kernel to
re-express. The shift-invert \emph{body} lifts cleanly to a global tensor-field
expression---it is just apply-$\Mop$, solve, untransform---but the \emph{loop}
around it is an opaque-library obstruction, a genuine boundary of what the
specification can rotate. The honest disposition is to document the boundary, not
to pretend a loop exists where the library hides one. The full constructive
realization of what the library does---Lanczos, Arnoldi, Krylov--Schur in our own
vocabulary---is \cref{ch:krylovschur}; here it stays a black box on purpose.
\end{pitfall}

\section{The quadratic eigenproblem: linearization}
\label{sec:quadratic}

The derivation so far assumed a lossless cavity, giving the linear pencil
$\Kop\vx=\lambda\Mop\vx$. Real cavities lose energy---to wall conductivity,
dielectric absorption, radiation through ports---and \cref{ch:reductions-pde}
showed that loss adds a damping operator $\Cop$ and makes the eigenproblem
\emph{quadratic} in the eigenvalue:
\begin{equation}
  \bigl(\Kop + \lambda\,\Cop + \lambda^2\,\Mop\bigr)\,\vx = \mathbf 0,
  \qquad \lambda = i\omega.
  \label{eq:qevp}
\end{equation}
Now $\lambda$ is complex: its imaginary part is the oscillation frequency, its
real part the decay rate. This is a \emph{polynomial eigenvalue problem}, and the
shift-invert iteration we built expects an ordinary or generalized \emph{linear}
problem. The standard remedy is \emph{linearization}: trade the quadratic
$N\times N$ problem for a linear $2N\times 2N$ one, at the cost of doubling the
size.

The trick is the same order-reduction we used for the transient ODE in
\cref{ch:reductions-pde}: introduce a second unknown for $\lambda\vx$. Set
$\vz=\lambda\vx$ and stack $(\vx,\vz)$. Then \cref{eq:qevp} becomes the linear
generalized eigenproblem
\begin{equation}
  \underbrace{\begin{pmatrix} \mathbf 0 & \mat{I} \\ -\Kop & -\Cop \end{pmatrix}}_{\tilde\Kop}
  \begin{pmatrix}\vx\\\vz\end{pmatrix}
  =
  \lambda
  \underbrace{\begin{pmatrix} \mat{I} & \mathbf 0 \\ \mathbf 0 & \Mop \end{pmatrix}}_{\tilde\Mop}
  \begin{pmatrix}\vx\\\vz\end{pmatrix},
  \label{eq:linearization}
\end{equation}
which one can multiply out to verify: the top block row reads $\vz=\lambda\vx$
(the definition of $\vz$), and the bottom row reads
$-\Kop\vx-\Cop\vz=\lambda\Mop\vz$, i.e.
$-\Kop\vx-\Cop\lambda\vx=\lambda^2\Mop\vx$, which rearranges to the quadratic
\cref{eq:qevp}. The doubled pencil $(\tilde\Kop,\tilde\Mop)$ is an ordinary
generalized eigenproblem of the kind \cref{sec:shift-invert} handles---so we apply
the \emph{same} shift-invert machinery to it, and the library solves it with the
\emph{same} black-box call, now on the $2N$-dimensional pencil. Each eigenpair of
the doubled problem yields the physical $(\lambda,\vx)$ directly; the extra block
$\vz=\lambda\vx$ is discarded.

For the still more general \emph{nonlinear} eigenproblem---where the operator
depends on $\lambda$ in a way no polynomial captures, $\bigl(\Kop+\lambda\Cop
+\lambda^2\Mop+\mat{A}_2(\lambda)\bigr)\vx=\mathbf 0$---linearization no longer
suffices, and one needs a genuinely nonlinear eigensolver built on
\emph{deflation} (finding eigenpairs one at a time and projecting out the ones
already found). That is the subject of \cref{ch:deflation}; we flag it here only
to place the quadratic case as the easy middle rung between the linear pencil and
the fully nonlinear problem.

\section{Reading the results}
\label{sec:results}

The library hands back a set of converged eigenpairs $(\lambda_i,\vx_i)$. Two
short steps turn them into the physical products the engineer asked for, and one
necessary cleanup keeps the answer honest.

\subsection{Eigenvalue to frequency, eigenvector to mode field}

The eigenvalue carries the frequency. For the lossless linear problem
$\lambda=\omega^2$, so
\begin{equation}
  \omega = \sqrt\lambda,\qquad
  f = \frac{\omega}{2\pi} = \frac{\sqrt\lambda}{2\pi}.
  \label{eq:freq}
\end{equation}
For the lossy quadratic problem $\lambda=i\omega$ is complex, so the frequency is
$\omega = \lambda/i = -i\lambda$, and the resonant frequency is read from its
imaginary part. The \emph{eigenvector} $\vx_i$ is, directly, the spatial pattern
of the mode field---the standing-wave shape the cavity rings in, the
three-dimensional analogue of the $\sin(n\pi x/L)$ modes we computed by hand in
\cref{ch:reductions-pde}.

The lossy case carries a bonus the lossless case cannot: the \emph{quality
factor} $Q$, the dimensionless measure of how slowly a mode loses energy. It
comes from the ratio of the imaginary part of $\lambda$ (the oscillation) to its
real part (the decay):
\begin{equation}
  Q = \frac{\abs{\Im\lambda}}{2\,\abs{\Re\lambda}}
    = \frac{\Im\omega}{2\,\Re\omega}\cdot(\text{convention}),
\end{equation}
a high $Q$ meaning a sharp, long-ringing resonance. The exact assembly of $Q$
from the complex eigenpair is an output reduction, treated with the other
postprocessing products in \cref{ch:reductions}; here we note only that the
eigenmode analysis is the one that \emph{produces} it.

\begin{figure}[t]
\centering
\begin{tikzpicture}[node distance=9mm]
  \node[data, minimum width=20mm] (pair) {$(\lambda_i,\vx_i)$};
  % eigenvalue branch
  \node[opbox, minimum width=26mm, right=12mm of pair, yshift=8mm] (lam)
    {$\lambda_i$ (eigenvalue)};
  \node[product, minimum width=30mm, right=12mm of lam] (freq)
    {$f_i=\sqrt{\lambda_i}/2\pi$};
  \node[opbox, minimum width=26mm, below=4mm of freq] (q)
    {$Q_i$ (from $\Im\lambda$)};
  % eigenvector branch
  \node[opbox, minimum width=26mm, right=12mm of pair, yshift=-8mm] (vec)
    {$\vx_i$ (eigenvector)};
  \node[product, minimum width=30mm, right=12mm of vec, yshift=-6mm] (mode)
    {mode field $\vE_i(\vx)$};
  \draw[flow] (pair.east) to[out=30,in=180] (lam.west);
  \draw[flow] (pair.east) to[out=-30,in=180] (vec.west);
  \draw[flow] (lam) -- (freq);
  \draw[flow] (lam.east) to[out=-20,in=160] (q.west);
  \draw[flow] (vec) -- (mode);
  \node[font=\scriptsize\itshape, simGray, below=1mm of mode]
    {(\cref{ch:reductions-pde} computed the 1-D version by hand)};
\end{tikzpicture}
\caption[From eigenpair to physical product]{Reading the eigensolver's output.
Each converged eigenpair $(\lambda_i,\vx_i)$ splits two ways. The
\emph{eigenvalue} $\lambda_i$ gives the resonant frequency $f_i=\sqrt{\lambda_i}/
2\pi$ (and, in the lossy case, the quality factor $Q_i$ from its imaginary part,
reduced in \cref{ch:reductions}). The \emph{eigenvector} $\vx_i$ \emph{is} the
mode field $\vE_i(\vx)$, the spatial standing-wave pattern---the
three-dimensional descendant of the hand-computed cavity modes of
\cref{ch:reductions-pde}.}
\label{fig:readout}
\end{figure}

\subsection{Spurious modes and the divergence-free projection}

There is one trap the raw eigensolver output sets, and the eigenmode analysis must
spring it before reporting results. The curl--curl operator $\Kop$ has a large
\emph{null space}---it annihilates every gradient field $\Grad\phi$, because
$\Curl\Grad\equiv\mathbf 0$ (the gauge freedom of \cref{ch:reductions-pde}). Every
such gradient field is an eigenvector with eigenvalue $\lambda=0$. These are not
physical resonances; they are \emph{spurious modes}, artifacts of the operator's
null space, clustered at and near zero frequency exactly where the eigensolver,
hunting the lowest eigenvalues, will dredge them up by the dozen.

Left in, they pollute the answer: the solver returns a heap of bogus zero-frequency
``modes'' that crowd out the real low resonances. The cure is the
\emph{divergence-free projection} of \cref{ch:bc}: project each candidate
eigenvector onto the subspace orthogonal to the gradients (the discrete Helmholtz
decomposition), removing its gradient component before it is accepted as a mode.
With the gradient null-space projected away, the spurious zero-frequency modes
vanish and the eigensolver's lowest eigenvalues are the genuine physical
resonances.

\begin{figure}[t]
\centering
\begin{tikzpicture}
\begin{axis}[
    name=raw,
    width=0.46\linewidth, height=38mm,
    title={\footnotesize raw eigensolver output},
    title style={font=\footnotesize, yshift=-1mm},
    xlabel={frequency $f$}, ylabel={},
    xmin=-0.3, xmax=5, ymin=0, ymax=1.4,
    axis lines=left, ytick=\empty, xtick={0,2,3.2,4.4},
    xticklabels={$0$,,,}, xticklabel style={font=\scriptsize},
    xlabel style={font=\scriptsize}, clip=false,
  ]
  % spurious cluster near zero (gray) + physical modes
  \addplot[ycomb, very thick, simGray] coordinates {
    (0,1.0)(0.12,0.9)(0.22,0.85)(0.05,0.95)(0.3,0.8)};
  \addplot[ycomb, very thick, simBlue, mark=*, mark size=1.6pt] coordinates {
    (2,1.0)(3.2,1.0)(4.4,1.0)};
  \node[font=\scriptsize\itshape, simGray, align=center] at (axis cs:0.6,1.3)
    {spurious\\(gradient null-space)};
\end{axis}
\begin{axis}[
    name=proj, at={(raw.east)}, anchor=west, xshift=10mm,
    width=0.46\linewidth, height=38mm,
    title={\footnotesize after divergence-free projection},
    title style={font=\footnotesize, yshift=-1mm},
    xlabel={frequency $f$}, ylabel={},
    xmin=-0.3, xmax=5, ymin=0, ymax=1.4,
    axis lines=left, ytick=\empty, xtick={0,2,3.2,4.4},
    xticklabels={$0$,,,}, xticklabel style={font=\scriptsize},
    xlabel style={font=\scriptsize}, clip=false,
  ]
  \addplot[ycomb, very thick, simBlue, mark=*, mark size=1.6pt] coordinates {
    (2,1.0)(3.2,1.0)(4.4,1.0)};
  \node[font=\scriptsize\itshape, simGreen, align=center] at (axis cs:0.9,1.3)
    {spurious removed};
  \node[font=\scriptsize, simBlue] at (axis cs:2,1.12) {$f_1$};
  \node[font=\scriptsize, simBlue] at (axis cs:3.2,1.12) {$f_2$};
  \node[font=\scriptsize, simBlue] at (axis cs:4.4,1.12) {$f_3$};
\end{axis}
\end{tikzpicture}
\caption[Spurious modes removed by the divergence-free projection]{The gradient
null space and its cure. \emph{Left:} the raw eigensolver output, hunting the
lowest eigenvalues, dredges up a cluster of \emph{spurious} zero-frequency modes
(gray)---gradient fields $\Grad\phi$ in the null space of the curl--curl operator
($\Curl\Grad\equiv\mathbf 0$, \cref{ch:reductions-pde})---crowding the genuine
physical resonances (blue). \emph{Right:} the divergence-free projection of
\cref{ch:bc} removes each eigenvector's gradient component, the spurious cluster
vanishes, and the lowest remaining eigenvalues are the true low resonances
$f_1,f_2,f_3$.}
\label{fig:spurious}
\end{figure}

\section{A worked power iteration}
\label{sec:worked-power}

We close the technical development by running power iteration on a small matrix
fully by hand, watching the iterate rotate toward the dominant eigenvector and the
Rayleigh quotient close in on the eigenvalue---the concrete realization of
\cref{sec:power} that the shift-invert machinery then redirects toward the low end.

\begin{workedexample}{Power iteration and the Rayleigh quotient}{power}
Take the symmetric matrix
\[
  \mat{A} = \begin{pmatrix} 2 & 1 \\ 1 & 2 \end{pmatrix},
\]
whose eigenvalues are $\mu_1=3$ (eigenvector $\vu_1=\tfrac{1}{\sqrt2}(1,1)$) and
$\mu_2=1$ (eigenvector $\vu_2=\tfrac{1}{\sqrt2}(1,-1)$). The dominant eigenpair is
$(3,\vu_1)$; the convergence ratio is $\abs{\mu_2/\mu_1}=1/3$. Start from
$\vx_0=(1,0)$---deliberately \emph{not} aligned with either eigenvector---and
iterate $\vx_{k+1}=\mat{A}\vx_k/\norm{\mat{A}\vx_k}$, recording the Rayleigh
quotient $\rho_k=\inner{\mat{A}\vx_k}{\vx_k}/\inner{\vx_k}{\vx_k}$ at each step.

\medskip
\emph{Step 0.} $\vx_0=(1,0)$;\;
$\mat{A}\vx_0=(2,1)$;\;
$\rho_0 = \dfrac{(2,1)\cdot(1,0)}{(1,0)\cdot(1,0)} = \dfrac{2}{1} = 2.000$.

\emph{Step 1.} Normalize: $\vx_1=(2,1)/\sqrt5=(0.894,0.447)$.\;
$\mat{A}\vx_1=(2.236,1.789)$;\;
$\rho_1 = \dfrac{(2.236)(0.894)+(1.789)(0.447)}{1} = 2.800$.

\emph{Step 2.} Normalize $\mat{A}\vx_1$: $\vx_2=(0.781,0.625)$.\;
$\mat{A}\vx_2=(2.187,2.031)$;\;
$\rho_2 = (2.187)(0.781)+(2.031)(0.625) = 2.976$.

\emph{Step 3.} $\vx_3=(0.732,0.681)$;\;
$\mat{A}\vx_3=(2.145,2.094)$;\;
$\rho_3 = (2.145)(0.732)+(2.094)(0.681) = 2.997$.

\emph{Step 4.} $\vx_4=(0.717,0.697)$;\;
$\rho_4 = 2.9997$.

\medskip
The iterate $\vx_k$ marches steadily from $(1,0)$ toward
$\vu_1\approx(0.707,0.707)$---the second component climbing $0\to0.447\to0.625\to
0.681\to0.697$ toward equality with the first---and the Rayleigh quotient closes
in on $\mu_1=3$: $2.000\to2.800\to2.976\to2.997\to2.9997$. The error in $\rho_k$
shrinks by roughly the factor $(\mu_2/\mu_1)^2=1/9$ per step (the \emph{squared}
ratio, the symmetric-case bonus of \cref{eq:rayleigh}): $1.0\to0.20\to0.024\to
0.003\to0.0003$, each about a ninth of the last. \Cref{fig:power-conv} plots the
convergence.

This is power iteration finding the \emph{dominant} eigenvalue $\mu_1=3$. To find
the \emph{smallest} eigenvalue $\mu_2=1$ instead---the analogue of the low cavity
resonance---we would run the same iteration on the shifted-inverted operator
$(\mat{A}-\sigma\mat{I})^{-1}$ with a shift $\sigma$ near $1$: by \cref{eq:eigmap}
the eigenvalue $\mu_2=1$ would become dominant in the transformed problem, and
this same iteration would converge to it instead.
\end{workedexample}

\begin{figure}[t]
\centering
\begin{tikzpicture}
\begin{axis}[
    width=0.66\linewidth, height=50mm,
    xlabel={iteration $k$}, ylabel={Rayleigh quotient $\rho_k$},
    xmin=0, xmax=4, ymin=1.8, ymax=3.15,
    axis lines=left, xtick={0,1,2,3,4}, ytick={2.0,2.5,3.0},
    xticklabel style={font=\scriptsize}, yticklabel style={font=\scriptsize},
    xlabel style={font=\scriptsize}, ylabel style={font=\scriptsize},
    legend style={font=\scriptsize, draw=simGray!50, at={(0.98,0.04)},
      anchor=south east},
    every axis plot/.append style={very thick},
  ]
  % the target eigenvalue
  \addplot[simGray, densely dashed, domain=0:4, samples=2] {3.0};
  \addlegendentry{$\mu_1=3$ (target)}
  % the Rayleigh quotient trajectory
  \addplot[simBlue, mark=*, mark size=2pt] coordinates {
    (0,2.000)(1,2.800)(2,2.976)(3,2.997)(4,2.9997)};
  \addlegendentry{$\rho_k$ (Rayleigh quotient)}
\end{axis}
\end{tikzpicture}
\caption[Power-iteration convergence of the Rayleigh quotient]{The Rayleigh
quotient $\rho_k$ of \cref{wex:power} converging to the dominant eigenvalue
$\mu_1=3$ (gray dashed). Starting from a generic vector, $\rho_k$ climbs
$2.000\to2.800\to2.976\to2.997\to2.9997$, its error falling by roughly the squared
ratio $(\mu_2/\mu_1)^2=1/9$ each step---the second-order accuracy the Rayleigh
quotient buys on a symmetric matrix. The fast, clean convergence here reflects a
well-separated spectrum ($\mu_1/\mu_2=3$); a clustered low spectrum would crawl,
which is exactly why the eigenmode problem needs the spectral transform.}
\label{fig:power-conv}
\end{figure}

\section{The shape of the eigenmode solve}

Step back and assemble the chapter into one picture. The eigenmode analysis,
alone among the six, solves an \emph{eigenvalue} problem rather than a linear
system or an evolution. It wants a few low resonances of a huge pencil
$(\Kop,\Mop)$---the hard, interior, lower end of the spectrum. The naive tool,
power/subspace iteration, finds the wrong (dominant) end and crawls on the
clustered low modes we care about. The spectral transform fixes both at once:
iterating against $(\Kop-\sigma\Mop)^{-1}\Mop$ remaps the eigenvalues so the ones
near the shift become dominant, at the cost of an inner linear solve per step. And
the whole outer iteration is delegated to a specialized library through reverse
communication, so Palace assembles the pencil, hands over the shift and the inner
solver, calls one opaque routine, and reads back the eigenpairs---which it turns
into frequencies, quality factors, and mode fields after projecting away the
spurious gradient null space.

\begin{keyidea}
The eigenmode solve is the deepest composition in the book: a
\emph{library-owned} outer eigenvalue iteration (\cref{ch:krylovschur}), calling
back through the \code{apply} interface (\cref{ch:operators}) into a Palace-owned
shift-invert body (apply $\Mop$, inner solve, untransform), whose inner solve is
itself a Palace-owned \code{ksp\_solve} (\cref{ch:krylov}) calling a
preconditioner (\cref{ch:multigrid}). Three nested ``a solver is an operator''
rotations, the outermost not even Palace's code. The spectral transform is the
hinge: it converts the hard interior eigenproblem into the easy dominant one the
iteration can solve, and pays for the conversion in inner linear solves.
\end{keyidea}

\summarysection
The eigenmode analysis seeks a few eigenpairs $(\lambda,\vx)$ of the generalized
eigenproblem $\Kop\vx=\lambda\Mop\vx$ ($\lambda=\omega^2$), with $\Kop$ the
curl--curl stiffness and $\Mop$ the mass---specifically the \emph{smallest}
nonzero $\lambda$, the lowest resonances at the hard interior/lower end of a vast
spectrum. \emph{Power iteration} $\vx\leftarrow\mat{A}\vx/\norm{\cdot}$ converges
to the \emph{dominant} (largest-magnitude) eigenvector at rate $\abs{\mu_2/\mu_1}$,
with the \emph{Rayleigh quotient} $\rho(\vx)=\inner{\mat{A}\vx}{\vx}/
\inner{\vx}{\vx}$ reading off the eigenvalue; subspace iteration finds several at
once. But this finds the wrong end and crawls on clustered spectra. The
\emph{spectral transform} (shift-and-invert) fixes both: iterating against
$\mat{T}_\sigma=(\Kop-\sigma\Mop)^{-1}\Mop$ maps each eigenvalue
$\lambda\mapsto\theta=1/(\lambda-\sigma)$, so the eigenvalues \emph{nearest the
shift} $\sigma$ become \emph{dominant}---turning ``find eigenvalues near $\sigma$''
(hard, interior) into ``find the dominant eigenvalues of $\mat{T}_\sigma$'' (easy).
The cost is an \emph{inner linear solve} of $(\Kop-\sigma\Mop)$ per apply---a
\code{ksp\_solve}, usually preconditioned---so the body is apply-$\Mop$
$\rhd$ inner-solve $\rhd$ untransform. In practice the \emph{outer} eigenvalue
iteration (Krylov--Schur/Arnoldi) lives inside a specialized library (SLEPc,
ARPACK) that drives the loop and \emph{calls back} for the apply (reverse
communication): the eigenmode driver assembles $(\Kop,\Mop)$, hands them plus the
shift to the library, and reads back converged eigenpairs in one opaque call---the
fourth, black-box solve corner. The lossy \emph{quadratic} eigenproblem
$(\Kop+\lambda\Cop+\lambda^2\Mop)\vx=\mathbf 0$ is \emph{linearized} to a doubled
linear pencil and solved the same way. Finally, eigenvalues become frequencies
$f=\sqrt\lambda/2\pi$ (and quality factors $Q$ from the imaginary part) and
eigenvectors become mode fields, after a \emph{divergence-free projection} removes
the spurious gradient null-space modes clustered at zero frequency.

\furtherreading
The definitive modern treatment of large-scale eigenvalue computation---power and
subspace iteration, the spectral transform, and the projection methods of the next
chapter---is Saad's \emph{Numerical Methods for Large Eigenvalue
Problems}~\citep{saad2011eig}; its chapters on the shift-and-invert Arnoldi method
are the direct background for \cref{sec:shift-invert,sec:blackbox}. The
Krylov--Schur restarting that the library uses to drive the outer iteration is due
to Stewart~\citep{stewart2001krylovschur}, the subject of \cref{ch:krylovschur}.
For the underlying linear-algebra fundamentals---the Rayleigh quotient, its
second-order accuracy, and the geometry of convergence---Trefethen and
Bau~\citep{trefethenbau1997} remain the clearest undergraduate-accessible source.
The reverse-communication black-box structure of \cref{sec:blackbox} is the
ARPACK/SLEPc design; the constructive realization of the loop those libraries hide
is developed, in our own vocabulary, in \cref{ch:krylovschur}, and the nonlinear
(deflation) generalization in \cref{ch:deflation}.

\exercisessection
\begin{enumerate}[nosep]
  \item \textbf{Which end, and why.} The eigenmode analysis wants the
  \emph{smallest} nonzero $\lambda$, yet power iteration finds the
  \emph{largest}. Using the eigenbasis expansion \cref{eq:power-expand}, explain
  in one sentence why iteration amplifies the dominant eigenvector. Then state the
  convergence ratio for the dominant pair and explain why a \emph{clustered} low
  spectrum (the cavity's near-degenerate overtones) makes \emph{direct} iteration
  for those modes especially slow.
  \item \textbf{The eigenvalue map by hand.} A pencil has eigenvalues
  $\lambda\in\{2,5,11,20\}$ and you want the two nearest $\lambda=6$. (i) Choose a
  shift $\sigma=6$ and compute the four transformed eigenvalues
  $\theta_i=1/(\lambda_i-\sigma)$. (ii) Rank them by magnitude and confirm the two
  you wanted are now dominant. (iii) The iteration returns $\theta=-1.0$; recover
  the physical $\lambda$. (iv) Why does the eigenvector need no untransforming?
  \item \textbf{Locating the cost.} Write out the three moves of one application of
  $\mat{T}_\sigma=(\Kop-\sigma\Mop)^{-1}\Mop$ to a vector $\vv$
  (\cref{eq:siapply}). Which move dominates the cost, and which chapter's machinery
  performs it? Explain why the shifted matrix $\Kop-\sigma\Mop$ is solved
  \emph{iteratively} rather than factored once, given that the shift $\sigma$ is
  held fixed for the whole eigensolve. (Hint: contrast with the fixed-operator
  electrostatic family of \cref{ch:reductions-pde}---when \emph{is} a single
  factorization the right call?)
  \item \textbf{Rayleigh quotient accuracy.} Run two more steps of the power
  iteration of \cref{wex:power} starting from $\vx_4=(0.717,0.697)$, computing
  $\rho_5$ and $\rho_6$. Verify the error in $\rho_k$ shrinks by roughly $1/9$ per
  step, and explain where the factor $1/9$ comes from. What would the per-step
  error-reduction factor be if instead $\mu_1=3$, $\mu_2=2.9$ (a near-cluster)?
  \item \textbf{The black box, drawn.} In your own words, describe \emph{reverse
  communication}: what does the library own, what does Palace own, and what is the
  single interface across which they cooperate? Why can the library not simply hold
  Palace's operators directly (\cref{ch:matrixfree})? Contrast the eigenmode solve
  corner with the electrostatic one: in which does Palace write the outer loop, and
  in which does it write only the apply?
  \item \textbf{Linearize the quadratic.} Verify by direct multiplication that the
  doubled pencil \cref{eq:linearization} is equivalent to the quadratic
  eigenproblem \cref{eq:qevp}: expand both block rows and recover
  $(\Kop+\lambda\Cop+\lambda^2\Mop)\vx=\mathbf 0$. What is the size of the doubled
  problem relative to the original, and what is discarded from each computed
  eigenvector when reading back the physical mode?
  \item \textbf{Spurious modes.} Explain why the curl--curl stiffness $\Kop$ has a
  large null space (cite the exactness identity from \cref{ch:reductions-pde}) and
  why every null-space vector is an eigenvalue-zero eigenvector of the pencil. Why
  do these spurious modes appear precisely where the eigenmode solver is looking
  (the low end)? What operation (\cref{ch:bc}) removes them, and what subspace does
  it project onto?
  \item \textbf{Reading the results.} A lossy eigensolve returns the complex
  eigenvalue $\lambda = i\,\omega$ with $\lambda = (-0.02 + 6.28\,i)\times10^{9}$.
  (i) What resonant frequency $f$ (in Hz) does this mode have? (ii) Which part of
  $\lambda$ carries the decay, and does a \emph{smaller} real part mean a higher or
  lower quality factor $Q$? (iii) Why does the lossless linear problem
  ($\lambda=\omega^2$ real) produce no $Q$ at all?
  \item \textbf{Choosing the shift.} You want the lowest five resonances of a
  cavity, and you know the fundamental is near $f_0$. (i) What shift $\sigma$ (in
  terms of $\lambda=\omega^2=(2\pi f)^2$) would you pass to the eigensolver, and
  why place it slightly below the lowest wanted eigenvalue rather than at zero?
  (ii) What goes wrong if you set $\sigma$ exactly equal to an eigenvalue? (iii) If
  the five wanted modes span a wide frequency range, why might a single shift
  converge slowly on the highest of them, and what does that suggest about the
  trade-off in choosing $\sigma$?
\end{enumerate}
